% Source-aligned English for printed p.162.
\SourcePage{162}

\LayerTarget{tohoku:prop:3.4.2:p2}
\textit{exact sequence of sheaves of groups on $X$, with $F$ abelian. Let
$E$ be a fiber space over $X$ with structure sheaf $H$ [11, Chapter IV], and
let $F^E$ be the sheaf of groups associated with $E$ and with the actions of
$H$ on $F$ defined through the inner automorphisms of $G$, which act on the
invariant sheaf $F$. One can then define, functorially, a ``boundary
element'' $\partial E\in H^2(X,F^E)$ such that $\partial E=0$ if and only if
the class $c(E)$ of $E$ in $H^1(X,H)$ [11, Chapter V] belongs to the image of
$H^1(X,G)$.}

This statement simplifies when $F$ lies in the center of $G$, for then
$F^E=F$ no longer depends on the fiber space $E$, and we obtain:

\LayerTarget{tohoku:cor:3.4.2}
\textsc{Corollary.} --- \textit{Let
$e\longrightarrow F\longrightarrow G\longrightarrow H\longrightarrow e$
be an exact sequence of sheaves of groups on $X$, with $F$ contained in the
center of $G$. There is a functorial map
\[
\partial:H^1(X,H)\longrightarrow H^2(X,F)
\]
such that $\partial^{-1}(0)=\operatorname{Im}(H^1(X,G))$.}

\LayerTarget{tohoku:proof:3.4.2:p1}
\textsc{Proof of Proposition 3.4.2.} Embed every sheaf $M$ on $X$ into the
sheaf $\overline M$ whose set of sections over an open subset $U$ is

\LayerTarget{tohoku:construction:flabby-envelope}
\[
\prod_{x\in U}M(x).
\]
Thus $\overline M$ is flabby, and $M$ identifies with a subsheaf of
$\overline M$. If $M$ is a sheaf of groups, or of abelian groups, the same is
true of $\overline M$. Moreover, if $L$ is a flabby sheaf of groups, not
necessarily abelian, then $H^1(X,L)$ consists only of the neutral element. In
other words, every principal $L$-sheaf [11, Definition 3.4.2] has a section;
one readily constructs such a section by ``Zornification'' over the set of
sections constructed on open subsets of $X$.

\LayerTarget{tohoku:comparison:cech-derived-h1}
It follows that, for a sheaf $F$ of abelian groups, the group $H^1(X,F)$
defined in [11] by the Čech method agrees with the group defined axiomatically
in the present work, namely the first right satellite $S^1\Gamma$ of
$\Gamma$ on $C^X$. Returning to the hypotheses of Proposition 3.4.2, we have
a homomorphism of exact sequences:

\LayerTarget{tohoku:diagram:3.4.2:flabby-embedding}
\[
\begin{array}{ccccccccc}
e&\longrightarrow&F&\longrightarrow&G&\longrightarrow&H&\longrightarrow&e\\
&&\big\downarrow&&\big\downarrow&&\big\downarrow&&\\[0.45em]
e&\longrightarrow&\overline F&\longrightarrow&\overline G&\longrightarrow
&\overline H&\longrightarrow&e
\end{array}
\qquad \overline F\cap G=F
\]
The sheaves $G$ and $\overline F$ are subsheaves of groups of $\overline G$,
with $\overline F$ normal. Let $\overline F\mathbin{\cdot}G$ be the subsheaf
of groups of $\overline G$ generated by $G$ and $\overline F$, and put

\LayerTarget{tohoku:construction:3.4.2:P-Fprime}
\LayerTarget{tohoku:eq:3.4.2:quotients}
\[
P=\overline F G,\qquad F'=\overline F/F.
\]
There is an evident exact sequence of homomorphisms of sheaves of groups

\LayerTarget{tohoku:eq:3.4.2:P-extension}
\[
e\longrightarrow\overline F\longrightarrow P\longrightarrow H
\longrightarrow e.
\]

\LayerTarget{tohoku:claim:3.4.2:H1-bijection}
We claim that the corresponding map
$H^1(X,P)\longrightarrow H^1(X,H)$ is bijective. It is a monomorphism by the
exact cohomology sequence of [11, Proposition 5.6.2], since $H^1(X,L)$
consists only of the neutral element

\TohokuContinuation{163}
